\subsection*{Futures, \texttt{.data}, and \texttt{.span}}

The result of a DMA statement is a future-like handle:

\begin{lstlisting}[language=croqtile,caption={A future from a copy}]
f = dma.copy input => shared;
\end{lstlisting}

That handle gives you two important things:

\begin{itemize}[leftmargin=1.6em]
\item \texttt{f.data}: the moved tile, as a spanned tensor.
\item \texttt{f.span}: the shape of that tile.
\end{itemize}

Example:

\begin{lstlisting}[language=croqtile,caption={Reading data and span}]
f = dma.copy input => shared;
foreach idx in [f.span]
  output.at(idx) = f.data.at(idx);
\end{lstlisting}

Even synchronous copies are worth naming when you want to pass the moved tile to later code.

\begin{figure}[htbp]\centering
\includegraphics[width=\linewidth,height=0.56\textheight,keepaspectratio]{f04-fig_future_data_light.png}
\caption{The copy result exposes the destination tensor through .data and its shape through .span; asynchronous reads require completion first.}
\end{figure}
